 Capítulo  1
 Interfaces  do  sistema  operacional
 A  função  de  um  sistema  operacional  é  compartilhar  um  computador  entre  vários  programas  e  fornecer  um  conjunto  de  serviços  
mais  útil  do  que  o  hardware  suporta  sozinho.  Um  sistema  operacional  gerencia  e  abstrai  o  hardware  de  baixo  nível,  de  modo  que,  
por  exemplo,  um  processador  de  texto  não  precise  se  preocupar  com  o  tipo  de  hardware  de  disco  que  está  sendo  usado.  Um  
sistema  operacional  compartilha  o  hardware  entre  vários  programas  para  que  eles  sejam  executados  (ou  pareçam  ser  executados)  
ao  mesmo  tempo.  Por  fim,  os  sistemas  operacionais  fornecem  maneiras  controladas  para  os  programas  interagirem,  para  que  
possam  compartilhar  dados  ou  trabalhar  em  conjunto.
 Um  sistema  operacional  fornece  serviços  aos  programas  do  usuário  por  meio  de  uma  interface.  Projetar  uma  boa  interface  
acaba  sendo  difícil.  Por  um  lado,  gostaríamos  que  a  interface  fosse  simples  e  concisa,  pois  isso  facilita  a  implementação  correta.  
Por  outro  lado,  podemos  ser  tentados  a  oferecer  muitos  recursos  sofisticados  aos  aplicativos.  O  segredo  para  resolver  essa  
tensão  é  projetar  interfaces  que  dependam  de  alguns  mecanismos  que  possam  ser  combinados  para  fornecer  muita  generalidade.
 Este  livro  utiliza  um  único  sistema  operacional  como  exemplo  concreto  para  ilustrar  conceitos  de  sistema  operacional.  Esse  
sistema  operacional,  o  xv6,  fornece  as  interfaces  básicas  introduzidas  pelo  sistema  operacional  Unix  de  Ken  Thompson  e  Dennis  
Ritchie  [17],  além  de  imitar  o  design  interno  do  Unix.
 O  Unix  oferece  uma  interface  estreita  cujos  mecanismos  se  combinam  bem,  oferecendo  um  grau  surpreendente  de  generalidade.  
Essa  interface  tem  sido  tão  bem-sucedida  que  os  sistemas  operacionais  modernos  —  BSD,  Linux,  macOS,  Solaris  e  até  mesmo,  
em  menor  grau,  o  Microsoft  Windows  —  possuem  interfaces  semelhantes  às  do  Unix.  Entender  o  xv6  é  um  bom  começo  para  
entender  qualquer  um  desses  sistemas  e  muitos  outros.
 Como  mostra  a  Figura  1.1,  o  xv6  assume  a  forma  tradicional  de  um  kernel,  um  programa  especial  que  fornece  serviços  aos  
programas  em  execução.  Cada  programa  em  execução,  chamado  de  processo,  possui  memória  contendo  instruções,  dados  e  
uma  pilha.  As  instruções  implementam  a  computação  do  programa.  Os  dados  são  as  variáveis  sobre  as  quais  a  computação  atua.  
A  pilha  organiza  as  chamadas  de  procedimento  do  programa.
 Um  determinado  computador  normalmente  tem  muitos  processos,  mas  apenas  um  único  kernel.
 Quando  um  processo  precisa  invocar  um  serviço  do  kernel,  ele  invoca  uma  chamada  de  sistema,  uma  das  chamadas  na  
interface  do  sistema  operacional.  A  chamada  de  sistema  entra  no  kernel;  o  kernel  executa  o  serviço  e  retorna.  Assim,  um  processo  
alterna  entre  a  execução  no  espaço  do  usuário  e  no  espaço  do  kernel.
 Conforme  descrito  em  detalhes  nos  capítulos  subsequentes,  o  kernel  usa  os  mecanismos  de  proteção  de  hardware  
fornecidos  por  uma  CPU1  para  garantir  que  cada  processo  em  execução  no  espaço  do  usuário  possa  acessar  apenas
 1Este  texto  geralmente  se  refere  ao  elemento  de  hardware  que  executa  uma  computação  com  o  termo  CPU,  uma  sigla
 9
Machine Translated by Google
 usuário
 espaço
 núcleo
 espaço
 concha
 gato
 chamada  de  
sistema
 Núcleo
 Figura  1.1:  Um  kernel  e  dois  processos  de  usuário.
 sua  própria  memória.  O  kernel  executa  com  os  privilégios  de  hardware  necessários  para  implementar  essas  proteções;  os  programas  do  
usuário  são  executados  sem  esses  privilégios.  Quando  um  programa  do  usuário  invoca  uma  chamada  de  sistema,  o  hardware  aumenta  
o  nível  de  privilégio  e  começa  a  executar  uma  função  pré-definida  no  kernel.
 O  conjunto  de  chamadas  de  sistema  que  um  kernel  fornece  é  a  interface  que  os  programas  do  usuário  veem.  O  kernel  xv6  fornece  
um  subconjunto  dos  serviços  e  chamadas  de  sistema  que  os  kernels  Unix  tradicionalmente  oferecem.
 A  Figura  1.2  lista  todas  as  chamadas  de  sistema  do  xv6.
 O  restante  deste  capítulo  descreve  os  serviços  do  xv6  —  processos,  memória,  descritores  de  arquivo,  pipes  e  um  sistema  de  
arquivos  —  e  os  ilustra  com  trechos  de  código  e  discussões  sobre  como  o  shell,  a  interface  de  linha  de  comando  do  Unix,  os  utiliza.  O  
uso  de  chamadas  de  sistema  pelo  shell  ilustra  o  cuidado  com  que  foram  projetados.
 O  shell  é  um  programa  comum  que  lê  comandos  do  usuário  e  os  executa.  O  fato  de  o  shell  ser  um  programa  do  usuário,  e  não  parte  do  kernel,  ilustra  o  poder  da  
interface  de  chamada  de  sistema:  não  há  nada  de  especial  no  shell.  Isso  também  significa  que  o  shell  é  fácil  de  substituir;  como  resultado,  os  sistemas  Unix  modernos  
têm  uma  variedade  de  shells  para  escolher,  cada  um  com  sua  própria  interface  de  usuário  e  recursos  de  script.  O  shell  xv6  é  uma  implementação  simples  da  essência  do  
shell  Bourne  do  Unix.  Sua  implementação  pode  ser  encontrada  em  (user/sh.c:1).
 1.1  Processos  e  memória
 Um  processo  xv6  consiste  em  memória  de  espaço  do  usuário  (instruções,  dados  e  pilha)  e  estado  por  processo,  privado  do  kernel.  O  Xv6  
compartilha  o  tempo  dos  processos:  ele  alterna  de  forma  transparente  as  CPUs  disponíveis  entre  o  conjunto  de  processos  aguardando  
execução.  Quando  um  processo  não  está  em  execução,  o  xv6  salva  os  registradores  de  CPU  do  processo,  restaurando-os  na  próxima  
execução.  O  kernel  associa  um  identificador  de  processo,  ou  PID,  a  cada  processo.
 Um  processo  pode  criar  um  novo  processo  usando  a  chamada  de  sistema  fork .  Fork  fornece  ao  novo  processo  uma  cópia  exata  da  
memória  do  processo  que  fez  a  chamada,  tanto  instruções  quanto  dados.  Fork  retorna  tanto  no  processo  original  quanto  no  novo.  No  
processo  original,  fork  retorna  o  PID  do  novo  processo.  No  novo  processo,  fork  retorna  zero.  Os  processos  original  e  novo  são  
frequentemente  chamados  de  pai  e  filho.
 para  unidade  central  de  processamento.  Outra  documentação  (por  exemplo,  a  especificação  RISC-V)  também  utiliza  as  palavras  processador,  
núcleo  e  hart  em  vez  de  CPU.
 10
Machine Translated by Google
 Chamada  de  sistema
 int  fork()  int  
exit(int  status)  int  
wait(int  *status)  int  kill(int  
pid)  int  getpid()  int  
sleep(int  n)  int  
exec(char  *file,  char  
Descrição
 Crie  um  processo  e  retorne  o  PID  do  filho.
 Encerra  o  processo  atual;  status  reportado  para  wait().  Sem  retorno.
 Aguarde  a  saída  de  uma  criança;  saia  do  status  em  *status;  retorna  o  PID  da  criança.
 Encerra  o  PID  do  processo.  Retorna  0  ou  -1  em  caso  de  erro.
 Retorna  o  PID  do  processo  atual.
 Pausa  por  n  tiques  de  relógio.
 *argv[])  Carrega  um  arquivo  e  o  executa  com  argumentos;  retorna  somente  se  houver  erro.
 char  *sbrk(int  n)  int  
open(char  *file,  int  flags)  int  write(int  
fd,  char  *buf,  int  n)  Grava  n  bytes  de  buf  no  descritor  de  arquivo  fd;  retorna  n.
 Aumenta  a  memória  do  processo  em  n  bytes.  Retorna  o  início  da  nova  memória.
 Abre  um  arquivo;  sinalizadores  indicam  leitura/gravação;  retorna  um  fd  (descritor  de  arquivo).
 int  read(int  fd,  char  *buf,  int  n)  Lê  n  bytes  em  buf;  retorna  o  número  lido;  ou  0  se  for  o  fim  do  arquivo.
 int  fechar(int  fd)
 Libere  o  arquivo  aberto  fd.
 int  dup(int  fd)
 int  pipe(int  p[])
 int  chdir(char  *dir)
 int  mkdir(char  *dir)
 int  mknod(char  *arquivo,  int,  int)
 int  fstat(int  fd,  struct  stat  *st)
 Retorna  um  novo  descritor  de  arquivo  referindo-se  ao  mesmo  arquivo  que  fd.
 Crie  um  pipe,  coloque  descritores  de  arquivo  de  leitura/gravação  em  p[0]  e  p[1].
 Altera  o  diretório  atual.
 Crie  um  novo  diretório.
 Crie  um  arquivo  de  dispositivo.
 Coloque  informações  sobre  um  arquivo  aberto  em  *st.
 int  stat(char  *file,  struct  stat  *st)  Coloque  informações  sobre  um  arquivo  nomeado  em  *st.
 int  link(char  *file1,  char  *file2)  Cria  outro  nome  (file2)  para  o  arquivo  file1.
 int  unlink(char  *arquivo)
 Remover  um  arquivo.
 Figura  1.2:  Chamadas  de  sistema  Xv6.  Se  não  indicado  de  outra  forma,  essas  chamadas  retornam  0  para  nenhum  erro  e  -1  se
 há  um  erro.
 Por  exemplo,  considere  o  seguinte  fragmento  de  programa  escrito  na  linguagem  de  programação  C  [7]:
 int  pid  =  garfo();
 se(pid  >  0){
 printf("pai:  filho=%d\n",  pid);
 pid  =  esperar((int  *)  0);
 printf("criança  %d  terminou\n",  pid);
 }  senão  se(pid  ==  0){
 printf("filho:  saindo\n");
 saída(0);
 }  outro  {
 printf("erro  de  bifurcação\n");
 }
 A  chamada  de  sistema  exit  faz  com  que  o  processo  de  chamada  pare  de  executar  e  libere  recursos  como
 memória  e  arquivos  abertos.  Exit  recebe  um  argumento  de  status  inteiro,  convencionalmente  0  para  indicar  sucesso
 e  1  para  indicar  falha.  A  chamada  de  sistema  wait  retorna  o  PID  de  um  filho  saído  (ou  eliminado)  de
 11
Machine Translated by Google
 o  processo  atual  e  copia  o  status  de  saída  do  filho  para  o  endereço  passado  para  wait;  se  nenhum  dos  filhos  do  chamador  
tiver  saído,  wait  aguarda  que  um  o  faça.  Se  o  chamador  não  tiver  filhos,  wait  retorna  imediatamente  -1.  Se  o  pai  não  se  
importar  com  o  status  de  saída  de  um  filho,  ele  pode  passar  um  endereço  0  para  wait.
 No  exemplo,  as  linhas  de  saída
 pai:  filho=1234  filho:  saindo
 pode  sair  em  qualquer  ordem  (ou  até  mesmo  misturados),  dependendo  se  o  pai  ou  o  filho  chega  primeiro  à  sua  
chamada  printf .  Após  a  saída  do  filho,  a  espera  do  pai  retorna,  fazendo  com  que  o  pai  imprima
 pai:  filho  1234  está  pronto
 Embora  o  processo  filho  tenha  inicialmente  o  mesmo  conteúdo  de  memória  que  o  processo  pai,  ambos  estão  
executando  com  memória  e  registradores  separados:  alterar  uma  variável  em  um  não  afeta  o  outro.  Por  exemplo,  
quando  o  valor  de  retorno  de  wait  é  armazenado  em  pid  no  processo  pai,  isso  não  altera  a  variável  pid  no  processo  
filho.  O  valor  de  pid  no  processo  filho  ainda  será  zero.
 A  chamada  de  sistema  exec  substitui  a  memória  do  processo  que  fez  a  chamada  por  uma  nova  imagem  de  
memória  carregada  de  um  arquivo  armazenado  no  sistema  de  arquivos.  O  arquivo  deve  ter  um  formato  específico,  
que  especifica  qual  parte  do  arquivo  contém  instruções,  qual  parte  são  dados,  em  qual  instrução  iniciar,  etc.  O  Xv6  
usa  o  formato  ELF,  que  o  Capítulo  3  discute  em  mais  detalhes.  Normalmente,  o  arquivo  é  o  resultado  da  compilação  
do  código-fonte  de  um  programa.  Quando  exec  é  bem-sucedido,  ele  não  retorna  ao  programa  que  fez  a  chamada;  
em  vez  disso,  as  instruções  carregadas  do  arquivo  começam  a  ser  executadas  no  ponto  de  entrada  declarado  no  
cabeçalho  ELF.  exec  recebe  dois  argumentos:  o  nome  do  arquivo  que  contém  o  executável  e  uma  matriz  de  
argumentos  de  string.  Por  exemplo:
 char  *argv[3];
 argv[0]  =  "eco";  argv[1]  =  
"olá";  argv[2]  =  0;  exec("/bin/
 eco",  argv);  
printf("erro  de  execução\n");
 Este  fragmento  substitui  o  programa  chamador  por  uma  instância  do  programa /bin/echo  em  execução  com  a  lista  
de  argumentos  echo  hello.  A  maioria  dos  programas  ignora  o  primeiro  elemento  da  matriz  de  argumentos,  que  é  
convencionalmente  o  nome  do  programa.
 O  shell  xv6  usa  as  chamadas  acima  para  executar  programas  em  nome  dos  usuários.  A  estrutura  principal  do  
shell  é  simples;  veja  main  (user/sh.c:146).  O  loop  principal  lê  uma  linha  de  entrada  do  usuário  com  getcmd.  Em  
seguida,  chama  fork,  que  cria  uma  cópia  do  processo  do  shell.  O  processo  pai  chama  wait,  enquanto  o  processo  
filho  executa  o  comando.  Por  exemplo,  se  o  usuário  tivesse  digitado  "echo  hello"  no  shell,  runcmd  teria  sido  
chamado  com  "echo  hello"  como  argumento.  runcmd  (user/sh.c:55)  executa  o  comando  real.  Para  "echo  hello",  ele  
chamaria  exec  (user/sh.c:79).  Se  exec  for  bem-sucedido,  o  filho  executará  instruções  de  echo  em  vez  de  runcmd.  
Em  algum  momento,  echo  chamará  exit,  o  que  fará  com  que  o  pai  retorne  de  wait  em  main  (user/sh.c:146).
 12
Machine Translated by Google
 Você  pode  se  perguntar  por  que  fork  e  exec  não  são  combinados  em  uma  única  chamada;  veremos  mais  
adiante  que  o  shell  explora  essa  separação  em  sua  implementação  de  redirecionamento  de  E/S.  Para  evitar  o  
desperdício  de  criar  um  processo  duplicado  e  substituí-lo  imediatamente  (por  exec),  os  kernels  operacionais  
otimizam  a  implementação  de  fork  para  esse  caso  de  uso  usando  técnicas  de  memória  virtual,  como  cópia  na  
gravação  (consulte  a  Seção  4.6).
 O  Xv6  aloca  a  maior  parte  da  memória  do  espaço  do  usuário  implicitamente:  fork  aloca  a  memória  necessária  
para  a  cópia  filha  da  memória  do  pai,  e  exec  aloca  memória  suficiente  para  armazenar  o  arquivo  executável.  Um  
processo  que  precisa  de  mais  memória  em  tempo  de  execução  (talvez  para  malloc)  pode  chamar  sbrk(n)  para  
aumentar  sua  memória  de  dados  em  n  bytes;  sbrk  retorna  a  localização  da  nova  memória.
 1.2  E/S  e  descritores  de  arquivo
 Um  descritor  de  arquivo  é  um  pequeno  inteiro  que  representa  um  objeto  gerenciado  pelo  kernel  do  qual  um  
processo  pode  ler  ou  gravar.  Um  processo  pode  obter  um  descritor  de  arquivo  abrindo  um  arquivo,  diretório  ou  
dispositivo,  criando  um  pipe  ou  duplicando  um  descritor  existente.  Para  simplificar,  frequentemente  nos  referimos  
ao  objeto  ao  qual  um  descritor  de  arquivo  se  refere  como  "arquivo";  a  interface  do  descritor  de  arquivo  abstrai  
as  diferenças  entre  arquivos,  pipes  e  dispositivos,  fazendo  com  que  todos  pareçam  fluxos  de  bytes.  Nos  
referiremos  à  entrada  e  à  saída  como  E/S.
 Internamente,  o  kernel  xv6  usa  o  descritor  de  arquivo  como  um  índice  em  uma  tabela  por  processo,  de  
modo  que  cada  processo  tenha  um  espaço  privado  de  descritores  de  arquivo  começando  em  zero.  Por  
convenção,  um  processo  lê  do  descritor  de  arquivo  0  (entrada  padrão),  grava  a  saída  no  descritor  de  arquivo  1  
(saída  padrão)  e  grava  mensagens  de  erro  no  descritor  de  arquivo  2  (erro  padrão).  Como  veremos,  o  shell  
explora  a  convenção  para  implementar  redirecionamentos  de  E/S  e  pipelines.  O  shell  garante  que  sempre  tenha  
três  descritores  de  arquivo  abertos  (user/sh.c:152).  que  são,  por  padrão,  descritores  de  arquivo  para  o  console.
 O  sistema  de  leitura  e  gravação  chama  bytes  de  leitura  e  bytes  de  gravação  para  abrir  arquivos  nomeados  por  
descritores  de  arquivo.  A  chamada  read(fd,  buf,  n)  lê  no  máximo  n  bytes  do  descritor  de  arquivo  fd,  copia-os  para  buf  e  
retorna  o  número  de  bytes  lidos.  Cada  descritor  de  arquivo  que  se  refere  a  um  arquivo  possui  um  deslocamento  associado.  
read  lê  os  dados  do  deslocamento  do  arquivo  atual  e,  em  seguida,  avança  esse  deslocamento  pelo  número  de  bytes  lidos:  
uma  leitura  subsequente  retornará  os  bytes  seguintes  aos  retornados  pela  primeira  leitura.  Quando  não  há  mais  bytes  
para  ler,  read  retorna  zero  para  indicar  o  fim  do  arquivo.
 A  chamada  write(fd,  buf,  n)  grava  n  bytes  de  buf  no  descritor  de  arquivo  fd  e  retorna  o  número  de  bytes  gravados.  
Menos  de  n  bytes  são  gravados  somente  quando  ocorre  um  erro.  Assim  como  read,  write  grava  dados  no  deslocamento  
do  arquivo  atual  e,  em  seguida,  avança  esse  deslocamento  pelo  número  de  bytes  gravados:  cada  gravação  continua  de  
onde  a  anterior  parou.
 O  seguinte  fragmento  de  programa  (que  constitui  a  essência  do  programa  cat)  copia  dados  de  sua  entrada  
padrão  para  sua  saída  padrão.  Se  ocorrer  um  erro,  ele  grava  uma  mensagem  na  entrada  padrão.
 erro.
 char  buf[512];  int  n;
 para(;;){
 13
Machine Translated by Google
 n  =  ler(0,  buf,  tamanho  de  buf);  se(n  ==  0)  quebrar;  
se(n  <  0)
 { fprintf(2,  
"erro  de  
leitura\n");  sair(1);
 }  se(escrever(1,  buf,  n) !=  n){
 fprintf(2,  "erro  de  gravação\n");  exit(1);
 }
 }
 O  importante  a  ser  observado  no  fragmento  de  código  é  que  cat  não  sabe  se  está  lendo  de  um  arquivo,  console  ou  pipe.  Da  
mesma  forma,  cat  não  sabe  se  está  imprimindo  para  um  console,  um  arquivo  ou  qualquer  outro  lugar.  O  uso  de  descritores  
de  arquivo  e  a  convenção  de  que  o  descritor  de  arquivo  0  é  a  entrada  e  o  descritor  de  arquivo  1  é  a  saída  permitem  uma  
implementação  simples  de  cat.
 A  chamada  de  sistema  close  libera  um  descritor  de  arquivo,  tornando-o  livre  para  reutilização  em  futuras  chamadas  de  
sistema  open,  pipe  ou  dup  (veja  abaixo).  Um  descritor  de  arquivo  recém-alocado  é  sempre  o  descritor  não  utilizado  de  
menor  numeração  do  processo  atual.
 Descritores  de  arquivo  e  fork  interagem  para  facilitar  a  implementação  do  redirecionamento  de  E/S.  O  fork  copia  a  tabela  
de  descritores  de  arquivo  do  processo  pai  juntamente  com  sua  memória,  de  modo  que  o  processo  filho  comece  com  exatamente  
os  mesmos  arquivos  abertos  que  o  processo  pai.  A  chamada  de  sistema  exec  substitui  a  memória  do  processo  que  fez  a  
chamada,  mas  preserva  sua  tabela  de  arquivos.  Esse  comportamento  permite  que  o  shell  implemente  o  redirecionamento  de  E/
 S  por  meio  de  fork,  reabrindo  os  descritores  de  arquivo  escolhidos  no  processo  filho  e,  em  seguida,  chamando  exec  para  
executar  o  novo  programa.  Aqui  está  uma  versão  simplificada  do  código  que  um  shell  executa  para  o  comando  cat  <  input.txt:
 char  *argv[2];
 argv[0]  =  "gato";  argv[1]  =  
0;  if(fork()  ==  0)  
{ fechar(0);  abrir("input.txt",  
O_RDONLY);  
exec("gato",  argv);
 }
 Após  o  processo  filho  fechar  o  descritor  de  arquivo  0,  open  garante  que  usará  esse  descritor  de  arquivo  para  o  
arquivo  input.txt  recém-aberto:  0  será  o  menor  descritor  de  arquivo  disponível.  cat  então  executa  com  o  descritor  de  
arquivo  0  (entrada  padrão)  referindo-se  ao  arquivo  input.txt.  Os  descritores  de  arquivo  do  processo  pai  não  são  
alterados  por  esta  sequência,  pois  ela  modifica  apenas  os  descritores  do  processo  filho.
 O  código  para  redirecionamento  de  E/S  no  shell  xv6  funciona  exatamente  dessa  maneira  (user/sh.c:83).  Lembre-se  de  
que  neste  ponto  do  código  o  shell  já  bifurcou  o  shell  filho  e  que  runcmd  chamará  exec  para  carregar  o  novo  programa.
 O  segundo  argumento  para  open  consiste  em  um  conjunto  de  sinalizadores,  expressos  como  bits,  que  controlam  
o  que  open  faz.  Os  valores  possíveis  são  definidos  no  cabeçalho  de  controle  de  arquivo  (fcntl)  (kernel/fcntl.h:1-5):
 14
Machine Translated by Google
 O_RDONLY,  O_WRONLY,  O_RDWR,  O_CREATE  e  O_TRUNC,  que  instruem  open  a  abrir  o  arquivo  para  leitura,  ou  para  
gravação,  ou  para  leitura  e  gravação,  para  criar  o  arquivo  se  ele  não  existir  e  para  truncá-lo  para  comprimento  zero.
 Agora  deve  ficar  claro  por  que  é  útil  que  fork  e  exec  sejam  chamadas  separadas:  entre  as  duas,  o  shell  tem  a  chance  de  
redirecionar  a  E/S  do  shell  filho  sem  perturbar  a  configuração  de  E/S  do  shell  principal.  Pode-se,  em  vez  disso,  imaginar  uma  
hipotética  chamada  de  sistema  combinada  forkexec ,  mas  as  opções  para  realizar  o  redirecionamento  de  E/S  com  tal  
chamada  parecem  estranhas.  O  shell  poderia  modificar  sua  própria  configuração  de  E/S  antes  de  chamar  forkexec  (e  então  
desfazer  essas  modificações);  ou  forkexec  poderia  receber  instruções  de  redirecionamento  de  E/S  como  argumentos;  ou  (o  
que  é  menos  atraente)  todo  programa,  como  cat,  poderia  ser  ensinado  a  fazer  seu  próprio  redirecionamento  de  E/S.
 Embora  o  fork  copie  a  tabela  do  descritor  de  arquivo,  cada  deslocamento  de  arquivo  subjacente  é  compartilhado  entre
 Pais  e  filhos.  Considere  este  exemplo:
 if(fork()  ==  0)  { write(1,  "olá  
",  6);  exit(0); }  else  { wait(0);  write(1,  
"mundo\n",  
6);
 }
 No  final  deste  fragmento,  o  arquivo  anexado  ao  descritor  de  arquivo  1  conterá  os  dados  hello  world.
 A  gravação  no  pai  (que,  graças  à  espera,  só  é  executada  após  a  conclusão  do  filho)  continua  de  onde  a  gravação  do  filho  
parou.  Esse  comportamento  ajuda  a  produzir  uma  saída  sequencial  a  partir  de  sequências  de  comandos  de  shell,  como  (echo  
hello;  echo  world)  >output.txt.
 A  chamada  de  sistema  dup  duplica  um  descritor  de  arquivo  existente,  retornando  um  novo  que  se  refere  ao  mesmo  
objeto  de  E/S  subjacente.  Ambos  os  descritores  de  arquivo  compartilham  um  deslocamento,  assim  como  os  descritores  de  
arquivo  duplicados  por  fork .  Esta  é  outra  maneira  de  escrever  hello  world  em  um  arquivo:
 fd  =  dup(1);  
escrever(1,  "olá  ",  6);  escrever(fd,  
"mundo\n",  6);
 Dois  descritores  de  arquivo  compartilham  um  deslocamento  se  forem  derivados  do  mesmo  descritor  de  arquivo  original  
por  uma  sequência  de  chamadas  fork  e  dup .  Caso  contrário,  os  descritores  de  arquivo  não  compartilham  deslocamentos,  
mesmo  que  tenham  resultado  de  chamadas  open  para  o  mesmo  arquivo.  O  dup  permite  que  shells  implementem  comandos  
como  este:  ls  existing-file  non-existing-file  >  tmp1  2>&1.  O  2>&1  informa  ao  shell  para  atribuir  ao  comando  um  descritor  de  
arquivo  2,  que  é  uma  duplicata  do  descritor  1.  Tanto  o  nome  do  arquivo  existente  quanto  a  mensagem  de  erro  para  o  arquivo  
inexistente  serão  exibidos  no  arquivo  tmp1.  O  shell  xv6  não  suporta  redirecionamento  de  E/S  para  o  descritor  de  arquivo  de  
erro,  mas  agora  você  sabe  como  implementá-lo.
 Os  descritores  de  arquivo  são  uma  abstração  poderosa,  porque  eles  escondem  os  detalhes  do  que  estão  conectados:  
um  processo  escrevendo  no  descritor  de  arquivo  1  pode  estar  escrevendo  em  um  arquivo,  em  um  dispositivo  como  o  console  
ou  em  um  pipe.
 15
Machine Translated by Google
 1.3  Tubos
 Um  pipe  é  um  pequeno  buffer  de  kernel  exposto  aos  processos  como  um  par  de  descritores  de  arquivo,  um  para  leitura  e  
outro  para  escrita.  A  gravação  de  dados  em  uma  extremidade  do  pipe  torna  esses  dados  disponíveis  para  leitura  na  outra  
extremidade.  Pipes  fornecem  uma  maneira  para  os  processos  se  comunicarem.
 O  código  de  exemplo  a  seguir  executa  o  programa  wc  com  a  entrada  padrão  conectada  à  extremidade  de  leitura  de  um  
pipe.
 int  p[2];  char  
*argv[2];
 argv[0]  =  "wc";  argv[1]  =  
0;
 pipe(p);  
if(fork()  ==  0)  { fechar(0);  
dup(p[0]);  
fechar(p[0]);  
fechar(p[1]);  exec("/
 bin/wc",  argv); }  else  
{ fechar(p[0]);  escrever(p[1],  "olá  
mundo\n",  
12);  fechar(p[1]);
 }
 O  programa  chama  pipe,  que  cria  um  novo  pipe  e  registra  os  descritores  de  arquivo  de  leitura  e  gravação  no  array  p.  Após  a  
bifurcação,  tanto  o  pai  quanto  o  filho  possuem  descritores  de  arquivo  referentes  ao  pipe.  O  filho  chama  close  e  dup  para  fazer  
com  que  o  descritor  de  arquivo  zero  se  refira  à  extremidade  de  leitura  do  pipe,  fecha  os  descritores  de  arquivo  em  p  e  chama  
exec  para  executar  wc.  Quando  wc  lê  de  sua  entrada  padrão,  ele  lê  do  pipe.  O  pai  fecha  o  lado  de  leitura  do  pipe,  escreve  no  
pipe  e,  em  seguida,  fecha  o  lado  de  gravação.
 Se  não  houver  dados  disponíveis,  uma  leitura  em  um  pipe  aguarda  que  os  dados  sejam  gravados  ou  que  todos  os  
descritores  de  arquivo  referentes  à  extremidade  de  gravação  sejam  fechados;  neste  último  caso,  a  leitura  retornará  0,  como  
se  o  fim  de  um  arquivo  de  dados  tivesse  sido  alcançado.  O  fato  de  a  leitura  bloquear  até  que  seja  impossível  a  chegada  de  
novos  dados  é  um  dos  motivos  pelos  quais  é  importante  que  o  filho  feche  a  extremidade  de  gravação  do  pipe  antes  de  
executar  o  comando  wc  acima:  se  um  dos  descritores  de  arquivo  de  wc  se  referisse  à  extremidade  de  gravação  do  pipe,  wc  
nunca  veria  o  fim  do  arquivo.
 O  shell  xv6  implementa  pipelines  como  grep  fork  sh.c  |  wc  -l  de  maneira  semelhante  ao  código  acima  (usuário/sh.c:101).  
O  processo  filho  cria  um  pipe  para  conectar  a  extremidade  esquerda  do  pipeline  com  a  extremidade  direita.  Em  seguida,  ele  
chama  fork  e  runcmd  para  a  extremidade  esquerda  do  pipeline  e  fork  e  runcmd  para  a  extremidade  direita,  e  aguarda  a  
conclusão  de  ambos.  A  extremidade  direita  do  pipeline  pode  ser  um  comando  que  inclui  um  pipe  (por  exemplo,  a  |  b  |  c),  que  
bifurca  dois  novos  processos  filhos  (um  para  b  e  outro  para  c).  Assim,  o  shell  pode  criar  uma  árvore  de  processos.  As  folhas
 16
Machine Translated by Google
 dessa  árvore  são  comandos  e  os  nós  internos  são  processos  que  esperam  até  que  os  filhos  esquerdo  e  direito  sejam  concluídos.
 Os  pipes  podem  não  parecer  mais  poderosos  do  que  os  arquivos  temporários:  o  pipeline
 eco  olá  mundo  |  wc
 poderia  ser  implementado  sem  tubos  como
 eco  olá  mundo  >/tmp/xyz;  wc  </tmp/xyz
 Pipes  têm  pelo  menos  três  vantagens  sobre  arquivos  temporários  nessa  situação.  Primeiro,  pipes  se  autolimpam;  com  o  
redirecionamento  de  arquivos,  um  shell  teria  que  ter  o  cuidado  de  remover /tmp/xyz  ao  terminar.  Segundo,  pipes  podem  passar  
fluxos  de  dados  arbitrariamente  longos,  enquanto  o  redirecionamento  de  arquivos  requer  espaço  livre  suficiente  em  disco  para  
armazenar  todos  os  dados.  Terceiro,  pipes  permitem  a  execução  paralela  de  etapas  do  pipeline,  enquanto  a  abordagem  de  
arquivo  exige  que  o  primeiro  programa  termine  antes  do  segundo  iniciar.
 1.4  Sistema  de  arquivos
 O  sistema  de  arquivos  xv6  fornece  arquivos  de  dados,  que  contêm  matrizes  de  bytes  não  interpretadas,  e  diretórios,  que  contêm  
referências  nomeadas  a  arquivos  de  dados  e  outros  diretórios.  Os  diretórios  formam  uma  árvore,  começando  em  um  diretório  
especial  chamado  raiz.  Um  caminho  como /a/b/c  refere-se  ao  arquivo  ou  diretório  chamado  c  dentro  do  diretório  chamado  b  
dentro  do  diretório  chamado  a  no  diretório  raiz /.  Caminhos  que  não  começam  com /  são  avaliados  em  relação  ao  diretório  atual  
do  processo  chamador,  que  pode  ser  alterado  com  a  chamada  de  sistema  chdir .  Ambos  os  fragmentos  de  código  abrem  o  mesmo  
arquivo  (assumindo  que  todos  os  diretórios  envolvidos  existam):
 chdir("/a");  chdir("b");  
open("c",  
O_RDONLY);
 aberto("/a/b/c",  O_RDONLY);
 O  primeiro  fragmento  altera  o  diretório  atual  do  processo  para /a/b;  o  segundo  não  faz  referência  nem  altera  o  diretório  atual  do  
processo.
 Existem  chamadas  de  sistema  para  criar  novos  arquivos  e  diretórios:  mkdir  cria  um  novo  diretório,  open  com  o  
sinalizador  O_CREATE  cria  um  novo  arquivo  de  dados  e  mknod  cria  um  novo  arquivo  de  dispositivo.  Este  exemplo  
ilustra  todos  os  três:
 mkdir("/dir");  fd  =  open("/
 dir/arquivo",  O_CREATE|O_WRONLY);  close(fd);  mknod("/console",  1,  
1);
 O  mknod  cria  um  arquivo  especial  que  se  refere  a  um  dispositivo.  Associados  a  um  arquivo  de  dispositivo  estão  os  números  de  
dispositivo  principal  e  secundário  (os  dois  argumentos  do  mknod),  que  identificam  exclusivamente  um  dispositivo  do  kernel.
 Quando  um  processo  abre  posteriormente  um  arquivo  de  dispositivo,  o  kernel  desvia  chamadas  de  leitura  e  gravação  do  sistema  
para  a  implementação  do  dispositivo  do  kernel  em  vez  de  passá-las  para  o  sistema  de  arquivos.
 17
Machine Translated by Google
 O  nome  de  um  arquivo  é  diferente  do  próprio  arquivo;  o  mesmo  arquivo  subjacente,  chamado  inode,  pode  ter  vários  
nomes,  chamados  links.  Cada  link  consiste  em  uma  entrada  em  um  diretório;  a  entrada  contém  um  nome  de  arquivo  e  
uma  referência  a  um  inode.  Um  inode  contém  metadados  sobre  um  arquivo,  incluindo  seu  tipo  (arquivo,  diretório  ou  
dispositivo),  seu  comprimento,  a  localização  do  conteúdo  do  arquivo  no  disco  e  o  número  de  links  para  um  arquivo.
 A  chamada  de  sistema  fstat  recupera  informações  do  inode  ao  qual  um  descritor  de  arquivo  se  refere.
 preenche  uma  estrutura  stat,  definida  em  stat.h  (kernel/stat.h)  como:
 #define  T_DIR
 T_FILE  #define  
T_DEVICE  3 //  Dispositivo
 estrutura  stat  {
 int  dev;
 1 //  Diretório
 2 //  Arquivo  #define  
//  Dispositivo  de  disco  do  sistema  de  arquivos  
//  Número  do  inode  uint  
ino;  short  
type; //  Tipo  de  arquivo  short  nlink; //  Número  
de  links  para  o  arquivo  uint64  size; //  Tamanho  do  arquivo  em  
bytes };
 A  chamada  do  sistema  de  link  cria  outro  nome  de  sistema  de  arquivo  que  se  refere  ao  mesmo  inode  que  um  existente
 arquivo  ing.  Este  fragmento  cria  um  novo  arquivo  chamado  a  e  b.
 abrir("a",  O_CREATE|O_WRONLY);  link("a",  "b");
 Ler  ou  escrever  em  a  é  o  mesmo  que  ler  ou  escrever  em  b.  Cada  inode  é  identificado  por  um  número  de  inode  exclusivo.  
Após  a  sequência  de  código  acima,  é  possível  determinar  que  a  e  b  se  referem  ao  mesmo  conteúdo  subjacente  inspecionando  
o  resultado  de  fstat:  ambos  retornarão  o  mesmo  número  de  inode  (ino)  e  a  contagem  de  nlinks  será  definida  como  2.
 A  chamada  de  sistema  unlink  remove  um  nome  do  sistema  de  arquivos.  O  inode  do  arquivo  e  o  espaço  em  disco  que  
contém  seu  conteúdo  só  são  liberados  quando  a  contagem  de  links  do  arquivo  é  zero  e  nenhum  descritor  de  arquivo  faz  
referência  a  ele.  Adicionando  assim
 desvincular("a");
 para  a  última  sequência  de  código  deixa  o  inode  e  o  conteúdo  do  arquivo  acessíveis  como  b.  Além  disso,
 fd  =  abrir("/tmp/xyz",  O_CREATE|O_RDWR);  desvincular("/tmp/
 xyz");
 é  uma  maneira  idiomática  de  criar  um  inode  temporário  sem  nome  que  será  limpo  quando  o  processo  fechar  o  fd  ou  sair.
 O  Unix  fornece  utilitários  de  arquivo  que  podem  ser  chamados  a  partir  do  shell  como  programas  de  nível  de  usuário,  por  
exemplo ,  mkdir,  ln  e  rm.  Esse  design  permite  que  qualquer  pessoa  estenda  a  interface  de  linha  de  comando  adicionando  
novos  programas  de  nível  de  usuário.  Em  retrospectiva,  esse  plano  parece  óbvio,  mas  outros  sistemas  projetados  na  época  
do  Unix  frequentemente  incorporavam  esses  comandos  ao  shell  (e  o  shell  ao  kernel).
 Uma  exceção  é  o  cd,  que  é  incorporado  ao  shell  (user/sh.c:161).  cd  deve  alterar  o  diretório  de  trabalho  atual  do  próprio  
shell.  Se  cd  fosse  executado  como  um  comando  regular,  o  shell
 18
Machine Translated by Google
 bifurcar  um  processo  filho,  o  processo  filho  executaria  cd,  e  cd  alteraria  o  diretório  de  trabalho  do  filho.  O  diretório  
de  trabalho  do  pai  (ou  seja,  do  shell)  não  mudaria.
 1.5  Mundo  real
 A  combinação  de  descritores  de  arquivo  "padrão",  pipes  e  sintaxe  shell  conveniente  do  Unix  para  operações  sobre  
eles  foi  um  grande  avanço  na  escrita  de  programas  reutilizáveis  de  uso  geral.  A  ideia  deu  início  a  uma  cultura  de  
"ferramentas  de  software"  que  foi  responsável  por  grande  parte  do  poder  e  popularidade  do  Unix,  e  o  shell  foi  a  
primeira  chamada  "linguagem  de  script".  A  interface  de  chamada  de  sistema  do  Unix  persiste  até  hoje  em  sistemas  
como  BSD,  Linux  e  macOS.
 A  interface  de  chamada  de  sistema  do  Unix  foi  padronizada  pelo  padrão  POSIX  (Portable  Operating  System  
Interface).  O  Xv6  não  é  compatível  com  POSIX:  faltam  muitas  chamadas  de  sistema  (incluindo  as  básicas,  como  
lseek),  e  muitas  das  chamadas  de  sistema  que  ele  fornece  diferem  do  padrão.  Nossos  principais  objetivos  para  o  
xv6  são  simplicidade  e  clareza,  ao  mesmo  tempo  em  que  fornecemos  uma  interface  de  chamada  de  sistema  
simples,  semelhante  ao  UNIX.  Várias  pessoas  estenderam  o  xv6  com  algumas  chamadas  de  sistema  a  mais  e  
uma  biblioteca  C  simples  para  executar  programas  Unix  básicos.  Os  kernels  modernos,  no  entanto,  fornecem  
muito  mais  chamadas  de  sistema  e  muitos  mais  tipos  de  serviços  de  kernel  do  que  o  xv6.  Por  exemplo,  eles  
oferecem  suporte  a  redes,  sistemas  de  janelas,  threads  em  nível  de  usuário,  drivers  para  muitos  dispositivos  e  
assim  por  diante.  Os  kernels  modernos  evoluem  contínua  e  rapidamente  e  oferecem  muitos  recursos  além  do  POSIX.
 O  Unix  unificou  o  acesso  a  múltiplos  tipos  de  recursos  (arquivos,  diretórios  e  dispositivos)  com  um  único  
conjunto  de  interfaces  de  nome  e  descritor  de  arquivo.  Essa  ideia  pode  ser  estendida  a  mais  tipos  de  recursos;  
um  bom  exemplo  é  o  Plan  9  [16],  que  aplicou  o  conceito  de  “recursos  são  arquivos”  a  redes,  gráficos  e  outros.  No  
entanto,  a  maioria  dos  sistemas  operacionais  derivados  do  Unix  não  seguiu  esse  caminho.
 O  sistema  de  arquivos  e  os  descritores  de  arquivos  têm  sido  abstrações  poderosas.  Mesmo  assim,  existem  
outros  modelos  para  interfaces  de  sistemas  operacionais.  O  Multics,  um  predecessor  do  Unix,  abstraiu  o  
armazenamento  de  arquivos  de  uma  forma  que  o  fez  parecer  memória,  produzindo  uma  interface  muito  diferente.  
A  complexidade  do  design  do  Multics  teve  uma  influência  direta  nos  designers  do  Unix,  que  buscavam  construir  
algo  mais  simples.
 O  Xv6  não  fornece  uma  noção  de  usuários  ou  de  proteção  de  um  usuário  de  outro;  em  termos  Unix,
 todos  os  processos  xv6  são  executados  como  root.
 Este  livro  examina  como  o  xv6  implementa  sua  interface  semelhante  à  do  Unix,  mas  as  ideias  e  conceitos  se  
aplicam  a  mais  do  que  apenas  o  Unix.  Qualquer  sistema  operacional  deve  multiplexar  processos  no  hardware  
subjacente,  isolar  processos  uns  dos  outros  e  fornecer  mecanismos  para  comunicação  controlada  entre  processos.  
Depois  de  estudar  o  xv6,  você  deverá  ser  capaz  de  analisar  outros  sistemas  operacionais  mais  complexos  e  
entender  os  conceitos  subjacentes  ao  xv6  nesses  sistemas  também
